\newpage
\section*{Day 3 Exercise 4 --- Blocking System }

In this exercise we simulate a blocking system (a loss system: $m$ service units, no waiting
room) by the \emph{event-by-event principle} and estimate the fraction of
blocked customers. 


The system has $m = 10$ servers and no waiting room: a customer that arrives
when all servers are busy is lost. With mean service time $s = 8$ and mean
inter-arrival time $1/\lambda = 1$, the offered traffic is $A = \lambda s = 8$
Erlang.

The simulation maps directly onto the four core elements of a discrete event
simulation (slides 5--6): the simulation clock is the time $t$ of the current
event, the state variable is the number of busy servers \texttt{n\_busy}, the
event list is a priority queue of arrival and departure events, and the
statistical accumulator is the \texttt{blocked} counter. The main loop advances
the clock to the next event, handles it, updates the state, and schedules
future events.

Each configuration is run as $10 \times 10{,}000$ customers. The 10 runs are
treated as independent \textbf{sub-samples} (slides 16--19): each yields one
blocked-fraction estimate $\hat\theta_i$, and a 95\% confidence interval is
formed from the $t$-distribution,
\begin{equation}
    \bar\theta \pm t_{0.975,\,9}\, \frac{S_\theta}{\sqrt{10}}.
    \label{eq:ci}
\end{equation}

The exact solution is Erlang's B-formula (Lecture 5),
\begin{equation}
    B = \frac{A^m / m!}{\sum_{i=0}^{m} A^i / i!},
    \label{eq:erlangb}
\end{equation}
which is \emph{insensitive to the service-time distribution} (only the mean
matters) but \emph{requires a Poisson arrival process}. For $A = 8$, $m = 10$
it gives $B = 0.1217$, the reference value throughout.

% ================================================================
\section*{Part 1 --- Verification (Poisson Arrivals, Exponential Service)}
% ================================================================

With Poisson arrivals (exponential inter-arrival times) and
exponential service, the simulation should reproduce Erlang B. The core event
loop accepts an arrival if a server is free, otherwise counts it as blocked:
\begin{verbatim}
while events:
    t, kind, i = heapq.heappop(events)        # advance clock to next event
    if kind == ARRIVAL:
        if n_busy < m:
            n_busy += 1
            heapq.heappush(events, (t + service[i], DEPARTURE, i))
        else:
            blocked += 1                       # all servers busy -> lost
    else:
        n_busy -= 1
\end{verbatim}
Erlang B (Equation~\ref{eq:erlangb}) is evaluated with a numerically stable
recursion that avoids large factorials:
\begin{verbatim}
b = 1.0
for i in range(1, m + 1):
    b = A * b / (i + A * b)
\end{verbatim}

\begin{table}[h]
    \centering
    \caption{Part 1 result, $10 \times 10{,}000$ customers. Erlang B $= 0.1217$.}
    \label{tab:part1}
    \renewcommand{\arraystretch}{1.1}
    \begin{tabular}{lrcc}
        \toprule
        Configuration & Fraction & 95\% CI & Erlang B in CI? \\
        \midrule
        Poisson / Exponential & $0.1221$ & $[0.1184,\ 0.1258]$ & Yes \\
        \bottomrule
    \end{tabular}
\end{table}

The simulated fraction $0.1221$ matches Erlang B and its CI contains the
analytical value, verifying the event-by-event implementation.

% ================================================================
\section*{Part 2 --- Renewal Arrival Processes (Exponential Service)}
% ================================================================

Replacing the Poisson process with a renewal process (general inter-arrival
times) keeps the mean arrival rate but changes its variability.
We test Erlang($2$) inter-arrivals and a hyperexponential mixture
($p_1 = 0.8, \lambda_1 = 0.8333,\ p_2 = 0.2, \lambda_2 = 5.0$), both with mean 1.

\begin{table}[h]
    \centering
    \caption{Part 2 results (exponential service). Erlang B $= 0.1217$.}
    \label{tab:part2}
    \renewcommand{\arraystretch}{1.1}
    \begin{tabular}{lrcc}
        \toprule
        Configuration & Fraction & 95\% CI & Erlang B in CI? \\
        \midrule
        Erlang($2$) arrivals      & $0.0950$ & $[0.0911,\ 0.0990]$ & No (lower)  \\
        Hyperexponential arrivals & $0.1398$ & $[0.1347,\ 0.1449]$ & No (higher) \\
        \bottomrule
    \end{tabular}
\end{table}

Erlang B no longer applies, since it assumes Poisson arrivals. Erlang($2$)
inter-arrival times have \emph{lower} variance than exponential (more regular,
less bursty), so customers are spaced more evenly, collide less, and blocking
\emph{drops} below Erlang B. Hyperexponential inter-arrival times have
\emph{higher} variance (a mix of very short and very long gaps), so customers
cluster and blocking \emph{rises} above Erlang B. Neither CI contains $0.1217$:
blocking depends on the \emph{variability} of the arrival process, not just its
mean rate.

% ================================================================
\section*{Part 3 --- Varied Service Distributions (Poisson Arrivals)}
% ================================================================

Returning to Poisson arrivals, we vary the service-time distribution while
keeping mean $8$. For the Pareto service times the scale is set so the mean
equals $8$, i.e.\ $\beta = 8(k-1)/k$:
\begin{verbatim}
beta = 8.0 * (k - 1) / k       # so that E[X] = beta*k/(k-1) = 8
\end{verbatim}

\begin{table}[h]
    \centering
    \caption{Part 3 results (Poisson arrivals). Erlang B $= 0.1217$.}
    \label{tab:part3}
    \renewcommand{\arraystretch}{1.1}
    \begin{tabular}{lrcc}
        \toprule
        Configuration & Fraction & 95\% CI & Erlang B in CI? \\
        \midrule
        Constant service   & $0.1202$ & $[0.1157,\ 0.1247]$ & Yes        \\
        Pareto $k=1.05$     & $0.0012$ & $[0.0002,\ 0.0022]$ & No (lower) \\
        Pareto $k=2.05$     & $0.1178$ & $[0.1126,\ 0.1229]$ & Yes        \\
        Lognormal service   & $0.1208$ & $[0.1160,\ 0.1257]$ & Yes        \\
        \bottomrule
    \end{tabular}
\end{table}

Erlang B is insensitive to the service-time distribution, confirmed by
constant, Pareto $k=2.05$, and lognormal service (all matching $0.1217$).
Pareto $k=1.05$ is the exception, but it is a \emph{finite-sample bias}, not a
real violation of insensitivity. With $k = 1.05$ the mean $E[X] = \beta k/(k-1)
= 8$ is finite, but the variance is infinite ($k < 2$). Almost every draw is
tiny ($\beta \approx 0.38$), and the rare enormous values needed to lift the
sample mean to $8$ essentially never appear in $10{,}000$ draws. The
\emph{empirical} mean service time is therefore far below $8$, the system is
under-loaded, and blocking collapses to $\approx 0$.

% ================================================================
\section*{Part 4 --- Comparison}
% ================================================================

Figure~\ref{fig:comparison} shows all configurations against the Erlang B reference line. \textbf{Constant, Pareto $k = 2.05$, and lognormal service} all have their CI crossing the line, confirming that Erlang B depends only on the mean
service time, not its shape. The Part 2 CIs miss the line entirely: \textbf{Erlang(2)} arrivals are more regular than Poisson, which means customers arrive evenly spaced, so fewer hit the system at once and blocking is lower. \textbf{Hyperexponential arrivals}
are more bursty, which means customers cluster together, filling all servers, and blocking is higher. The direction of the deviation matches the variability: less variable arrivals give less blocking, more variable give more. \textbf{Pareto
$k = 1.05$} sits far below the line with the narrowest CI of all (every run is biased in the same direction). The infinite-variance distribution cannot be sampled reliably with 10{,}000
observations, so the effective mean service time is far below 8 in every run and blocking is consistently near zero.

\begin{figure}[h]
    \centering
    \includegraphics[width=0.85\textwidth]{plots/day3/comparison.png}
    \caption{Blocked fraction with 95\% sub-sample CIs for all configurations
             against the Erlang B reference ($0.1217$, dashed). Green intervals
             contain Erlang B; red intervals deviate from it.}
    \label{fig:comparison}
\end{figure}
